\subsection{E.5. The two-parameter family and its actual singular fibres}
\begin{proposition}[E.5.1 (Refuted as stated; corrected): the input configuration]
The configuration in the question is not the configuration of
\cite{ViaTwo26push}.  The elliptic surface
$y^2=x^3+tx+1$ has fibre $III^*$ at infinity.  The constructed
threefold has two rank-one fibres at $p_1,p_2$, a rank-two
hexagonal fibre at $p_3$, and the multiple bielliptic fibre
$4B_*$ at infinity.
\end{proposition}
\begin{proof}
Put $s=t^{-1}$, $X=s^2x$, and $Y=s^3y$.  The minimal equation is
$Y^2=X^3+s^3X+s^6$, with coefficient and discriminant orders
$(3,6,9)$; this is the $III^*$ equation, as in
\cite[Lemma 2.1]{ViaTwo26push}.  The discriminant at finite
points is $-16(4t^3+27)$ with three simple zeros.  In the
threefold's multiplicative periods the orders at these points
are respectively $(1,0),(1,0),(1,1)$:
\[
 h(z,v)=(q\alpha(q)z,\beta(q)v),\qquad
 k(z,v)=(\gamma(q)z,\lambda q^\epsilon\delta(q)v),
\]
where all four coefficients are units, $\epsilon=0$ at $p_1,p_2$,
and $\epsilon=1$ at $p_3$.  Thus the first two torus ranks
drop by one and the last by two.  The explicit matrices of
E.5.2 independently have nilpotent ranks $1,1,2$; their exact
calculation rules out the assertion that all three are rank one.
These are the types stated in Proposition 3.1 and Theorem 1.1
of the cited source, not a correction to that source.
\end{proof}

\begin{proposition}[E.5.2 (Proved): monodromy, centralizer, and integral lattices]
Use the source's left-action marking
$H^1(F,\mathbb Z)=\mathbb Z\langle\psi,u,w,\delta\rangle$
and the ordered meridians $*,1,2,3$ with product one.
The cohomological matrices are
\[
\begin{aligned}
 T_*&=\begin{pmatrix}1&0&-2&-1\\0&0&-1&0\\0&1&0&-1\\0&0&0&1\end{pmatrix},&
 T_1&=\begin{pmatrix}1&0&2&2\\0&1&1&1\\0&0&1&0\\0&0&0&1\end{pmatrix},\\
 T_2&=\begin{pmatrix}1&0&0&0\\0&1&0&0\\0&-1&1&0\\0&0&0&1\end{pmatrix},&
 T_3&=\begin{pmatrix}1&0&0&-1\\0&1&1&0\\0&0&1&0\\0&0&0&1\end{pmatrix}.
\end{aligned}
\]
On ordinary dual homology put $H_i=(T_i^{-1})^t$,
$\ell=(1,-1,-1,0)^t$ and $\widehat\delta=(0,0,0,1)^t$.
Then
\[
 \operatorname{Cent}_{M_4(\mathbb Q)}(H_*,H_1,H_2,H_3)
 =\mathbb Q I\oplus\mathbb Q E_{41},
\]
and the cohomological centralizer is
$\mathbb Q I\oplus\mathbb Q E_{14}$.
The fixed and norm lattices at the multiple fibre are
\[
 \begin{aligned}
 \Lambda^{H_*}&=\mathbb Z\ell\oplus\mathbb Z\widehat\delta,\\
 N_4\Lambda&=4\mathbb Z\ell\oplus2\mathbb Z\widehat\delta,\\
 \Lambda^{H_*}/N_4\Lambda&\simeq\mathbb Z/4\oplus\mathbb Z/2.
 \end{aligned}
\]
The primitive vanishing lattices at the three reduced fibres are
\[
 K_1=\mathbb Z(\widehat w+\widehat\delta),\qquad
 K_2=\mathbb Z\widehat u,\qquad
 K_3=\mathbb Z\widehat w\oplus\mathbb Z\widehat\delta.
\]
Their invariant lattices are
\[
 \begin{split}
 \Lambda^{H_1}&=\mathbb Z(\widehat\psi-2\widehat u)
       \oplus\mathbb Z\widehat w\oplus\mathbb Z\widehat\delta,\\
 \Lambda^{H_2}&=\mathbb Z\widehat\psi\oplus
       \mathbb Z\widehat u\oplus\mathbb Z\widehat\delta,\\
 \Lambda^{H_3}&=K_3.
 \end{split}
\]
\end{proposition}
\begin{proof}
The matrices are Proposition 4.2 of \cite{ViaTwo26push}, whose
derivation uses the elliptic monodromies, the winding of $P$,
and the primitive rank-two cusp.  In particular it determines
the initially undetermined winding integer by the saturation
index $d^2=1$, rather than assuming this winding from the
desired topology.  Direct multiplication gives
$T_*T_1T_2T_3=I$, $T_*^4=I$ and
$(T_i-I)^2=0$ for $i=1,2,3$.
For clarity the homological matrices used in our calculation are
\[
\begin{aligned}
 H_*&=\begin{pmatrix}1&0&0&0\\-2&0&-1&0\\0&1&0&0\\1&1&0&1\end{pmatrix},&
 H_1&=\begin{pmatrix}1&0&0&0\\0&1&0&0\\-2&-1&1&0\\-2&-1&0&1\end{pmatrix},\\
 H_2&=\begin{pmatrix}1&0&0&0\\0&1&1&0\\0&0&1&0\\0&0&0&1\end{pmatrix},&
 H_3&=\begin{pmatrix}1&0&0&0\\0&1&0&0\\0&-1&1&0\\1&0&0&1\end{pmatrix}.
\end{aligned}
\]
Solving the sixty-four rational equations $CH_i=H_iC$
gives $C=aI+bE_{41}$ with free $a,b$; the retained exact
row-reduction program displays a basis and verifies its
dimension two.  Inverse transpose gives the stated
cohomological centralizer.  Solving $(H_*-I)v=0$ gives
$v=(a,-a,-a,d)^t$ integrally, and
\[
 N_4=I+H_*+H_*^2+H_*^3
 =\begin{pmatrix}4&0&0&0\\-4&0&0&0\\-4&0&0&0\\0&2&-2&4\end{pmatrix}.
\]
In the basis $(\ell,\widehat\delta)$ its columns are
$(4,0),(0,2),(0,-2),(0,4)$, giving exactly the displayed
norm group.  Each $K_i$ is the actual integral image of
$H_i-I$; its Smith form has only unit nonzero entries.
The kernels of these three matrices give the remaining fixed
lattices.  All products, rational row reductions, Hermite forms,
and Smith forms are executed in
\texttt{code/e\_push\_two\_parameter.py}.
\end{proof}

We now specify the parameter domain, because the source proves
a family on a sufficiently small bidisc, not on the entire
$\mu$-plane.  Set
\[
 S_\mu:\ y^2=x^3+tx+\mu t+1,\qquad
 P_\mu=(-\mu,\sqrt{1-\mu^3}),
\]
with the square root equal to one at zero.  Label the three roots
$p_i(\mu)$ of $4t^3+27(\mu t+1)^2$ holomorphically near zero;
$p_3$ is the zero of the chosen linearization.
Choose a holomorphic generator of its one-dimensional space of
sections vanishing at $p_3$ and scale it by $\lambda$.
For the added homology generator take the inverse of the
contracting lift.  On the universal cover of the four-punctured
base its period columns then have the form
\[
 \Pi_{\mu,c}=
 \begin{pmatrix}2z_\mu&\tau_\mu&1&0\\
 \beta_\mu+c&z_\mu&0&1\end{pmatrix},
 \qquad \lambda=e^{-2\pi i c},\qquad
 \partial_c\Pi=e\otimes\psi.
\]
Here $z_\mu$ is a lifted elliptic logarithm in the chosen marking;
changing the generator of the linearization changes
$\beta_\mu$ and the origin of $c$ by a holomorphic function of
$\mu$, and changes neither the integer-difference criterion
below nor the original $\lambda$-labelled bidisc.
This choice of inverse lift fixes the sign of $c$.
At $p_3$ its second period has jump $-1$, exactly canceled by
$H_3\widehat\psi=\widehat\psi+\widehat\delta$.
The covariance identities for all four matrices, including
$\tau_*=-1/\tau$, $z_*=(z-1)/\tau$,
$\beta_* =\beta-1-2(z-1)^2/\tau$, are checked symbolically
in the same program.

\begin{theorem}[E.5.3 (Proved): classification on the specified parameter domain]
There are $\epsilon>0$ and a real constant $C$ such that the
source construction, with its fixed logarithmic vector $\ell$,
defines a proper holomorphic family
$X_{\mu,c}$ over
$\{ |\mu|<\epsilon\}\times\{\operatorname{Im}c<C\}$.
It has the singular fibres of E.5.1 and the markings above.
Two of its members are biholomorphic if and only if
\[
 \mu'=\mu,\qquad c'-c\in2\mathbb Z.
\]
Every such biholomorphism covers the identity of the elliptic
base.  The quotient parameter in the second variable is
$v=e^{-\pi i c}$, with $\lambda=v^2$.
On a sufficiently small bidisc centred at any
$(0,\lambda_0)$ in Theorem 1.1 of \cite{ViaTwo26push}, choose
one logarithm of $\lambda$.  Within that actual bidisc,
\[
 X_{\mu,\lambda}\simeq X_{\mu',\lambda'}
 \quad\Longleftrightarrow\quad
 (\mu,\lambda)=(\mu',\lambda').
\]
\end{theorem}
\begin{proof}
The equations of Proposition 6.4 in \cite{ViaTwo26push}
give the simultaneous elliptic family, line bundle, linearization,
and good-reduction comparison.  On a smaller closed $\mu$-disc
all unit factors in its four local fillings are uniformly bounded
above and below.  Its contraction estimates therefore have a
common nonzero bound $|\lambda|<\eta$.  Pulling back by
$\lambda=e^{-2\pi i c}$ gives the indicated half-plane.
The integral markings and the four-torsion section $\ell/4$
continue uniquely on this simply connected parameter domain.
On every relatively compact subset the source's finite root
quotients and the fixed rank-one and hexagonal fans give
holomorphic charts with uniform bounds.  Their comparisons are
the specified section-preserving ones, so they agree on
intersections and define a family.  Properness is local on
the parameter domain and follows from the compact four-piece
construction there; smoothness follows from the same free
finite action and unimodular charts.  We make no assertion
about continuation to arbitrary $\mu$.

The intrinsic fibration is the anticanonical pencil.  To verify
the needed input from Proposition 6.1 of the source, on the
smooth core the line-coordinate form $dv/v$ wedged with a
base two-form identifies $K_X$ with $f^*\mathcal O(-1)$.
The forms $dq\wedge d\log z\wedge d\log v$ extend without
zeros or poles in both rank-one charts $q=ab$ and rank-two
charts $q=abc$.  At the finite root at infinity the source's
coordinate change gives
\[
 ds\wedge\frac{dX}{Y}\wedge\frac{dv}{v}
 =4\,dr\wedge\frac{dx'}{y'}\wedge\frac{dv'}{v'},
 \qquad s=r^4.
\]
It is invariant under the actual affine deck map and descends
through its free quotient.  Thus the same canonical formula
holds everywhere.  Connected fibres and projection formula
give $H^0(-K_X)=f^*H^0(\mathcal O(1))$, without base points.
A biholomorphism therefore induces an affine automorphism of
the elliptic base, fixing its distinguished multiple point
at infinity and its distinguished hexagonal point $p_3$.
It may initially interchange the two rank-one points.

Put $Q(\mu)=(p_1-p_3)/(p_2-p_3)$.
The source computes $p_i'(0)=-9/(2p_i(0))$.
If $Q'(0)$ vanished, the three velocities would have the form
$Ap_i(0)+B$; the polynomial $Az^2+Bz+9/2$ would vanish
at three distinct numbers, a contradiction.
Also $Q(0)^2-Q(0)+1=0$, so $Q(0)\ne Q(0)^{-1}$.
Shrink $\epsilon$ so that $Q$ is injective and its image
is disjoint from the image of $Q^{-1}$.  Preservation or
interchange of the rank-one values gives respectively
$Q(\mu')=Q(\mu)$ or $Q(\mu')=Q(\mu)^{-1}$.
The second possibility is excluded, and the first gives
$\mu'=\mu$.  The base automorphism now fixes three distinct
points individually and is the identity.

On regular fibres a biholomorphism is affine, and its integral
linear map commutes with all four $H_i$.  E.5.2 gives
$aI+bE_{41}$.  Comparing the constant third and fourth
period columns makes its complex linear part $aI$;
unimodularity gives $a=\alpha\in\{1,-1\}$.
The first column then gives $h=c'-c\in\mathbb Z$ and
lattice map $\alpha(I-hE_{41})$.
Normalize after the actual fourth-root base change.  Its
affine generator is the one with fourth power $\ell$, not
an arbitrary generator of its cyclic subgroup.  Summing the
translation defect around this generator gives
\[
 (I+hE_{41})\ell-\alpha\ell
 =(1-\alpha)\ell+h\widehat\delta\in N_4\Lambda.
\]
The norm lattice in E.5.2 forces $\alpha=1$ and $h$ even.
This proves necessity, including the multiple fibre's full
non-reduced structure.

We prove sufficiency with all four collar conditions.
Let $h$ be even and put
\[
 b_*=-\tfrac h4(\widehat u+\widehat w),\qquad
 L_*=-\tfrac h2\widehat u.
\]
Direct multiplication gives
\[
 (I-H_*)b_*=\tfrac h4\widehat\delta+L_*.
\]
Consequently $\zeta\mapsto\zeta+\Pi_{\mu,c}b_*$ on the
complete root torus family conjugates the two affine actions,
modulo the integral arrow $L_*$.  It and its inverse extend
over $r=0$ and descend.  On the punctured root collar its
middle translation includes
$h(\log r)/(2\pi i)e$; this is the derivative-free collar
term for the change of the invariant translation
$\Pi\ell/4$, not a logarithm used on the central chart.

At a rank-one point the vanishing coordinate is the first
toric coordinate in its chain cover, while the other
direction is the universal cover of the elliptic factor.
An extending translation has integral winding only in $K_j$.
After removing the corresponding power of $q$, its toric
coefficient is a holomorphic unit and its elliptic translation
extends holomorphically.  These are the full analytic collar
conditions.  Necessity follows by lifting a transverse disc
through a smooth central point to this free cover: its residual
defect in $\Lambda/K_j$ is zero.  The coordinate ratios on
the cover then give the asserted unit and holomorphic
remainder.  Conversely an integral power of $q$ translates
the infinite chain by an integer number of components;
the holomorphic remainders give torus and elliptic translations.
These maps and their inverses extend at both branches of every
node.  At the hexagonal point the same description uses two
toric windings in $K_3$, and integer translation of the fan
extends over its three double curves and both triple points.
This criterion includes regularity of the remainder; membership
of the defect in $K_j$ alone would not exclude an essential
singularity of a preassigned punctured translation.

Choose the three integral defects
\[
 L_1=0,\qquad L_2=\tfrac h2\widehat u,\qquad
 L_3=-\tfrac h2\widehat w.
\]
They belong to $K_1,K_2,K_3$, respectively, and are realized
by the preceding toric translations with initially constant
unit remainder.  In primitive cocharacter coordinates these
are $1$, $q_2^{h/2}$, and $(q_3^{-h/2},1)$, where the
monodromy defect of $q^a$ is defined as $a$ times the chosen
vanishing generator.  The required sphere relation is
\[
 L_*+H_*L_1+H_*H_1L_2+H_*H_1H_2L_3=0.
\]
It holds as an exact integer identity, checked in the program.
Thus the four local choices have compatible additive defects
and define a torsor of allowed middle translations.

We verify the vanishing of its analytic obstruction, including
the difference sheaf.  The root curve has just the order-four
point at infinity and is $\mathfrak C=\mathbb P(1,4)$,
with $K_{\mathfrak C}=\mathcal O(-5)$ and coarse pullback
$\mathcal O_B(1)=\mathcal O_{\mathfrak C}(4)$.
Vertical logarithmic fields on a rank-one chart are generated
by $a\partial_a-b\partial_b$ and the elliptic vector;
on a hexagonal chart they are the two height-zero toric fields.
The root chart has the ordinary vertical tangent frame.
Their linear transition functions depend only on the base;
they define a rank-two bundle $\mathcal E$ on $\mathfrak C$.
The fixed line-coordinate field supplies a subline
$\mathcal O\subset\mathcal E$ everywhere in these logarithmic
frames.  The logarithmic determinant sequence and cancellation
of the three reduced fibre divisors give
$K_X=\bar f^*(K_{\mathfrak C}\otimes(\det\mathcal E)^{-1})$.
The verified canonical formula and
$\bar f_*\mathcal O_X=\mathcal O_{\mathfrak C}$ therefore
give $\det\mathcal E=\mathcal O(-1)$.
The extension splits since
$H^1(\mathfrak C,\mathcal O(1))=0$, by the same two-chart
Laurent calculation with weights $1,4$.
The coarse difference sheaf is $V_B=\rho_*\mathcal E$:
on the root chart this is precisely finite equivariance,
and on each reduced collar it is precisely the holomorphic
remainder just described.  Finite averaging is exact, so
\[
 H^1(B,V_B)=H^1(\mathfrak C,\mathcal E)
 =H^1(\mathcal O\oplus\mathcal O(-1))=0.
\]
The torsor's Cech class vanishes.  Correcting the four local
translations by its holomorphic zero-cochain gives a global
middle translation and exact equality with all four local
extensions.  Their inverses have the opposite allowed
translations and glue as well.  This constructs the required
biholomorphism over $B$.

Finally choose a single logarithm of $\lambda$ on the source's
small bidisc.  Its $c$-values can be placed in a strip of
width less than two.  The criterion just proved then says
$c'=c$ and $\mu'=\mu$, which is equivalent to equality of
the original two parameters.  Conversely equality gives the
identity, with the fixed root, rank-one and hexagonal charts.
\end{proof}

The matrix data and the independent-variation argument are in
\cite[Proposition 4.2 and Proposition 6.4]{ViaTwo26push}; the
source states dimension at least two and does not state the
integer classification on the logarithmic lift.  We derived
the latter here from the marked norm and collar calculations.
For E.1--E.5 our searches on 15 November 2026 included the
exact arXiv identifiers together with ``classification'',
``Kuranishi'', ``deformation'', and ``parameters'', and the
references and version-specific records of the attached
papers.  No further paper resolving the general deformation
equivalence in E.2.4 was located.  This is a record of the
search performed, not a completeness or priority assertion;
the query responses and opened-source metadata are retained
in \texttt{evidence/E\_search*.json}.
